\section{Study 1: 2D Experiments}
\label{sec:study1}

Across three 2D datasets, two clear winners emerge. \textbf{SB-CFM} achieves the best endpoint quality on all three datasets (SWD@64 and MMD@64). \textbf{OT-CFM} produces the straightest trajectories (lowest curvature) on all three datasets and wins low-NFE robustness on two of three. At NFE${\leq}8$, OT and SB variants achieve SWD ${\sim}0.022$, while VP and Target reach ${\sim}0.040$---roughly 47\% worse. Table~\ref{tab:wins} summarizes the multi-dataset comparison.

\begin{table}[H]
\centering
\caption{Win counts across 3 datasets and 5 metrics (out of 3 possible wins each).}
\label{tab:wins}
\small
\begin{tabular}{@{}lccccc@{}}
\toprule
Variant & SWD@64 & MMD@64 & Low-NFE & Curvature & AUC \\
\midrule
OT-CFM & 0 & 0 & \textbf{2} & \textbf{3} & 0 \\
VP-CFM & 0 & 0 & 0 & 0 & 0 \\
Target-CFM & 0 & 0 & 0 & 0 & 1 \\
SB-CFM & \textbf{3} & \textbf{3} & 1 & 0 & \textbf{2} \\
\bottomrule
\end{tabular}
\end{table}

Crucially, \textbf{all four variants work flawlessly in 2D}---no instabilities, no divergence, smooth convergence across all datasets. This will prove misleading: the 2D setting masks numerical issues that become catastrophic in high dimensions (Section~\ref{sec:discussion}).
